\documentclass[12pt, a4paper] {article}
\usepackage{graphicx} % Required for inserting images
\usepackage{amsmath}
\usepackage{lmodern}
\usepackage{booktabs}
\usepackage{anyfontsize}
\graphicspath{{images/}}




%\title {\Huge\toprule Doppio pendolo con base oscillante \bottomrule}
%\author{\\Vittorio \textsc{Latrofa}\\Marta \textsc{Rosi} \\\\ }
%\date{Da inserire}

\begin{document}

%\maketitle

\newpage

\section{Descrizione del sistema}
Il sistema è composto da un carrello oscillante, al cui centro è saldata l'estremità superiore di un doppio pendolo. Il doppio pendolo è a sua volta un sottosistema composto da due aste, collegate in serie l'una all'altra. Il centro del carrello nell'istante iniziale è posizionato nell'origine del piano cartesiano scelto come riferimento. Ogni asta ha la sua lunghezza $l_i$ e massa tracurabile. All'estremità di ognuna di queste vi è una massa, rispettivamente $m_i$.
Sia infine $\theta_i$ l'angolo fra l'asse verticale e l'i-esima asta del pendolo.
\\\\
immagine
\section {Calcolo Lagrangiana}
L'evoluzione temporale del sistema verrà indagata attraverso il calcolo della Lagrangiana e la sua derivazione. In particolare si avranno, nel caso del doppio pendolo, 2 gradi di libertà dettati dagli angoli $\theta_1$ e $\theta_2$, e quindi il problema si ridurrà alla risoluzione di un sistema in due equazioni di Eulero-Lagrange.
\subsection{Posizioni e velocità cartesiane}
Per arrivare a definire la Lagrangiana necessitiamo dell'energia cinetica T e dell'energia potenziale U del sistema. Ricaviamo quindi posizioni e velocità cartesiane delle masse, a partire dagli angoli $\theta_1$ e $\theta_2$.\\
Partendo dalla prima massa, possiamo esprimere la sua posizione in coordinate cartesiane $(x_1,y_1)$ così:
\[ x_1=x_0+l_1sin(\theta_1)\hspace{1cm}y_1=-l_1cos(\theta_1)\]
Ponendo come $x_0$ l'ascissa del centro del carrello, che si sposta con legge oraria costante pari a:\[x_0=Asin(\omega t)\]
È interessante notare come esclusivamente la componente orizziontale è influenzata dalla forzatura in ingresso.\\\\
Proseguendo con la seconda massa, la sua posizione sarà descritta da:\[x_2=x_0+l_1sin(\theta_1)+l_2sin(\theta_2)\hspace{1cm}y_2=-l_1cos(\theta_1)-l_2cos(\theta_2)\]

(Nota: Ho fatto i calcoli anche per N-dimensioni, nel caso si volesse ampliare, però evito di scriveli perché è uno sbatti)

Deriviamo ora per ottenere le velocità:\[\dot x_n=A\omega cos(\omega t)+\sum_{i=1}^{n} l_i\dot\theta_i cos(\theta_i)\hspace{1cm}\dot y_n=\sum_{i=1}^{n}l_i\dot\theta_i sin(\theta_i)\hspace{1cm}n\in[1,2]\]
\subsection{Calcolo energia cinetica}
Calcoliamo ora l'energia cinetica del sistema. Consideriamo prima il contributo di una singola asta: \[T_i=\frac{m_i}{2}(\dot x_i^2+\dot y_i^2)\]
Calcoliamoci separatamente $\dot x_i^2$ e $\dot y_i^2$. Sapendo che \[(\sum_{k=1}^{n}a_k)^2=\sum_{i=1}^{n}\sum_{j=1}^{n}a_ia_j\]
possiamo scrivere (con $n\in[1,2]$):
\[\dot y_n^2=(\sum_{i=1}^{n}l_i\dot\theta_i sin(\theta_i))^2=\sum_{k=1}^{n}\sum_{j=1}^{n}l_kl_k\dot\theta_k \dot\theta_j sin(\theta_k)sin(\theta_j)\]
\[\dot x_n^2=(A\omega cos(\omega t)+\sum_{i=1}^{n}l_i\dot\theta_i cos(\theta_i))^2=\]
\[=A^2\omega ^2 cos^2 (\omega t)+2A\omega cos(\omega t)\sum_{i=1}^{n}l_i\dot\theta_icos(\theta_i)+\sum_{k=1}^{n}\sum_{j=1}^{n}l_kl_k\dot\theta_k \dot\theta_j cos(\theta_k)cos(\theta_j)\]
\\Sostituiamo 1 e 2 ad n, per ottenere $T_1$ e $T_2$, e quindi l'energia cinetica totale pari a:
\[T=T_1+T_2=\]
\[=\frac{M_{tot}}{2}(A^2\omega ^2cos^2(\omega t)+2A\omega cos(\omega t)l_1\dot \theta_1 cos(\theta_1)+l_1^2\dot \theta_1^2)+\]
\[+\frac{m_2}{2}(l_2^2\theta_2^2+2A\omega cos(\omega t)l_2\dot\theta_2cos(\theta_2)+2l_1l_2\dot\theta_1\dot\theta_2cos(\theta_1-\theta_2))\]\\
Con $M_{tot}=m_1+m_2$.
\subsection{Calcolo energia potenziale}
Ricaviamo l'energia potenziale U.\[U=U_1+U_2=m_1gy_1+m_2gy_2=-M_{tot}gl_1cos(\theta_1)-m_2gl_2cos(\theta_2)\]
\newpage
\subsection{Lagrangiana}
A questo punto possiamo esplicitare la Lagrangiana: \[L=T-U=\]\[=\frac{M_{tot}}{2}(A^2\omega ^2cos^2(\omega t)+2A\omega cos(\omega t)l_1\dot \theta_1 cos(\theta_1)+l_1^2\dot \theta_1^2+2gl_1cos(\theta_1))+\]\[+\frac{m_2}{2}(l_2^2\dot\theta_2^2+2A\omega cos(\omega t)l_2\dot\theta_2cos(\theta_2)+2l_1l_2\dot\theta_1\dot\theta_2cos(\theta_1-\theta_2)+2gl_2cos(\theta_2))\]
\section{Derivazione Lagrangiana}
Dalle equazioni di Eulero-Lagrange sappiamo che:
\[\left\{
\begin{aligned}
\frac {d}{dt}\frac{\partial L}{\partial \dot\theta_1}-\frac{\partial L}{\partial \theta_1}=0 \\
\frac {d}{dt}\frac{\partial L}{\partial \dot\theta_2}-\frac{\partial L}{\partial \theta_2}=0 
\end{aligned}
\right.\]
Procediamo calcolando le derivate parziali, partendo dalla prima equazione:
\[\frac{\partial L}{\partial \theta_1}=-M_{tot}(A\omega cos(\omega t)l_1\dot\theta_1 sin(\theta_1)+gl_1sin(\theta_1))-m_2(l_1l_2\dot\theta_1\dot\theta_2sin(\theta_1-\theta_2))\]
%%
\[\frac{\partial L}{\partial \dot\theta_1}=M_{tot}(l_1^2\dot \theta_1+A\omega cos(\omega t)l_1cos(\theta_1))+m_2(l_1l_2\dot\theta_2cos(\theta_1-\theta_2))\]
A questo punto deriviamo nel tempo:
\[\frac{d}{dt}\frac{\partial L}{\partial \dot\theta_1}=M_{tot}(l_1^2\ddot\theta_1-A\omega cos(\omega t)l_1\dot \theta_1 sin(\theta_1)-A\omega^2 sin(\omega t)l_1cos(\theta_1))+\]\[+m_2(l_1l_2\ddot\theta_2cos(\theta_1-\theta_2)-l_1l_2\dot\theta_2(\dot\theta_1-\dot\theta_2)sin(\theta_1-\theta_2))\]
Continuando con la seconda equazione:
\[\frac{\partial L}{\partial \theta_2}=m_2(l_1l_2\dot\theta_1\dot\theta_2sin(\theta_1-\theta_2)-A\omega cos(\omega t)l_2\dot\theta_2sin(\theta_2)-gl_2sin(\theta_2))\]
\[\frac{\partial L}{\partial \dot\theta_2}=m_2(l_2^2\dot\theta_2+A\omega cos(\omega t)l_2cos(\theta_2)+l_1l_2\dot\theta_1cos(\theta_1-\theta_2))\]
A questo punto derivimao nel tempo:
\[\frac{d}{dt}\frac{\partial L}{\partial \dot\theta_2}=m_2(l_2^2\ddot\theta_2-A\omega^2sin(\omega t)l_2cos(\theta_2)-A\omega cos(\omega t)l_2\dot\theta_2sin(\theta_2)+\]\[+l_1l_2\ddot\theta_1cos(\theta_1-\theta_2)-l_1l_2\dot\theta_1(\dot\theta_1-\dot\theta_2)sin(\theta_1-\theta_2))\]
\newpage
Uniamo quindi tutto e riscriviamo il sistema:
\[\
\begin{cases}
M_{tot}(l_1^2\ddot\theta_1-A\omega^2sin(\omega t)l_1cos(\theta_1)+gl_1sin(\theta_1))+m_2(l_1l_2\ddot\theta_2cos(\theta_1-\theta_2)+l_1l_2\dot\theta_2^2sin(\theta1-\theta_2)=0 \\
l_2^2\ddot\theta_2-A\omega^2 sin(\omega t)l_2cos(\theta_2)+l_1l_2\ddot\theta_1 cos(\theta_1-\theta_2)-l_1l_2\dot\theta_1^2sin(\theta_1-\theta_2)=0 
\end{cases}
\]


\end{document}
